\chapter{Results}
\label{chap:results}

{\small This chapter reports the main observations in four sections: object tracking quality (Section~\ref{sec:results_tracking}), runtime (Section~\ref{sec:results_runtime}), the comparison of VLMs and prompts (Section~\ref{sec:results_vlm}) and the resulting scene graphs across environments (Section~\ref{sec:results_scene_graph}). Each section states its objective and experimental setup before the results are discussed.\par}

\section{Object Tracking Quality}
\label{sec:results_tracking}

{\small The tracking evaluation addresses two questions: why object nodes occasionally contain two physical objects, and how reliably objects are re-identified when the pipeline is restarted in an environment it has already mapped. Both were examined on the uHumans2 office bag~\cite{rosinol2021kimera}, replayed at a reduced rate with ground-truth odometry on the Jetson AGX Orin, % TODO: verify hardware and playback rate of the tracking campaign
using the association of Section~\ref{sec:method_tracking}. Two instruments were used: a log with one record per mask-to-track decision, including the branch taken and the cue scores, and periodic per-track crops, saved every tenth observation, which show the raw single-frame mask and thus allow it to be distinguished from the accumulated extent of the track. All measurements come from a single environment, and since run lengths differ, track counts are normalised per 100 frames.\par}

{\small Inspection of the baseline configuration, which accepts a match when two of the five association verdicts agree and uses a centroid gate of 0.75\,m, showed that two mechanically distinct failure modes produce the same symptom, a single node whose box contains two objects. In \emph{segmentation-time fusion}, NanoSAM~\cite{nanosam2023} returns one mask spanning two touching surfaces, such as floor and sofa, wall and plant, or ceiling and pipe. Six or more instances were confirmed; the fusion persisted under every segmentation parameter tried and could begin mid-track, after about 50 clean observations, depending on the viewing angle. Since only one candidate reaches the tracker, no tracking parameter can repair it. In one floor track, both failure modes occurred about 20 observations apart.\par}

{\small In \emph{cross-track association merge}, a clean track absorbs a separately segmented object, as confirmed for a floor absorbing a wall and a sofa. For an object resting on a surface, the overlap, containment and height cues express one geometric fact: the height test passed on all 582 recorded associations, so the quorum of two was met before the centroid and image cues, the only ones able to object, had any influence. A single marginal acceptance (0.326 against a gate of 0.30) then dragged the floor track towards the sofa. Reweighting was not a remedy, since the centroid score of an erroneous merge (0.45) exceeded those of correct continuations (0.07, 0.24). The remedy therefore had to act on the acceptance condition itself.\par}

{\small Accordingly, the acceptance condition was tightened step by step over seven runs, each retaining the previous changes: a quorum of three with a centroid gate of 0.55\,m, a partial-overlap threshold of 0.50 with the centroid gate relaxed to 0.60\,m, the touch shortcut of the overlap test disabled, and finally the depth gate aligned with the minimum number of valid depth points required for geometry. The outcome is a deliberate trade. Raising the quorum increased fragmentation from 3.2 to between 8.0 and 9.9 tracks per 100 frames, and disabling the touch shortcut raised it to a maximum of 11.1. Part of this increase consisted of depth-less tracks, 16 of 101 in that run: masks of distant objects passed the original depth gate with only a few valid depth points, yielded no geometry and opened a new track at every observation. The aligned depth gate removed them entirely, leaving about 8.5 tracks per 100 frames, while the confirmed over-merges (floor and sofa, floor and trash can, wall and plant, ceiling and pipe) no longer occurred. This follows the asymmetric-error principle of Section~\ref{sec:method_tracking}: fragments can be merged later, whereas an over-merged node has wrong geometry, crop and label. Segmentation-time fusion remains unresolved (Chapter~\ref{chap:limitations}).\par}

{\small With this final configuration, re-identification was tested in five consecutive sessions, each replaying the same 300\,s wall-clock segment at a playback rate of 0.1, corresponding to 147 to 148 processed frames, with session persistence enabled (Section~\ref{subsec:method_sessions}). Memory was cleared only before session~1; session persistence is disabled in all other experiments of this chapter. In session~1, all 53 objects necessarily opened new tracks. Across sessions 2 to 5, only 1, 3, 0 and 1 new tracks were created, so that 199 of 204 re-encountered objects (97.5\,\%) were re-identified and the total number of tracks grew only from 53 to 58. Three of the five misses lie within about 1\,m of each other above 3.3\,m height, suggesting one tall structure whose centroid shifts beyond the gate between viewpoints. VLM calls fell monotonically from 28 to 11, 5, 2 and finally 1, since restored tracks keep their labels, while throughput stayed at about 0.49 frames per second (Section~\ref{sec:results_runtime}).\par}